% Part: lambda-calculus
% Chapter: introduction
% Section: primitive-recursion

\documentclass[../../../include/open-logic-section]{subfiles}

\begin{document}

\olfileid[id]{lam}{int}{pr}
\olsection{Fungsi yang \usetoken{S}{lambda definable} Tertutup terhadap Rekursi Primitif}

Untuk rekursi primitif, akhirnya kita perlu melakukan pekerjaan yang lebih
substantif. Kita harus melangkah secara bertahap. Seperti sebelumnya, dengan
mengandaikan bahwa kita telah mempunyai suku $G$ dan $H$ yang masing-masing
!!{lambda define} fungsi $g$ dan~$h$, kita ingin memperoleh suku $F$ yang
!!{lambda define}s fungsi~$f$ yang didefinisikan oleh
\begin{align*}
f(0, \vec z) & = g(\vec z) \\
f(x+1, \vec z) & = h(x, f(x,\vec z), \vec z).
\end{align*}
Jadi, secara umum, jika diberikan suku lambda $G'$ dan~$H'$, cukup ditemukan
suatu suku~$F$ sedemikian sehingga
\begin{align*}
F(\num{0}, \vec z) & \equiv G(\vec z) \\
F(\overline{n+1}, \vec z) & \equiv H(\num{n}, F(\num{n}, \vec z), \vec z)
\end{align*}
untuk setiap bilangan asli~$n$; fakta bahwa $G'$ dan~$H'$ !!{lambda define}
$g$ dan~$h$ berarti bahwa setiap kali kita memasukkan numeral
$\num{\vec m}$ sebagai pengganti $\vec z$, $F(\num{n+1}, \num{\vec m})$
akan ternormalisasi ke jawaban yang benar.

Untuk tujuan tersebut, cukup ditemukan suatu suku~$F$ yang memenuhi
\begin{align*}
F(\num 0) & \equiv G \\
F(\overline {n+1}) & \equiv H(\num{n},F(\num{n}))
\intertext{untuk setiap bilangan asli $n$, dengan}
G & = \lambd[\vec z][G'(\vec z)] \text{ dan}\\
H(u,v) & = \lambd[\vec z][H'(u,v(\vec z),\vec z)].
\end{align*}
Dengan kata lain, melalui kiat lambda kita tidak perlu lagi menangani parameter
tambahan $\vec z$ secara terpisah---parameter itu terserap ke dalam notasi
lambda.

Sebelum mendefinisikan suku $F$, kita memerlukan mekanisme untuk menangani
pasangan terurut. Mekanisme itu diberikan oleh lema berikut.

\begin{lem}
Terdapat suatu suku lambda $D$ sedemikian sehingga, bagi setiap pasangan suku
lambda $M$ dan~$N$, $D(M,N)(\num{0}) \red M$ dan
$D(M,N)(\num{1}) \red N$.
\end{lem}

\begin{proof}
Pertama, definisikan suku lambda $K$ dengan
\[
K(y) = \lambd[x][y].
\]
Dengan kata lain, $K$ adalah suku $\lambd[y][\lambd[x][y]]$. Dilihat dengan
cara lain, untuk setiap $M$, $K(M)$ adalah fungsi konstan yang mengembalikan~$M$
pada setiap masukan.

Sekarang definisikan $D(x,y,z)$ dengan $D(x,y,z)=z(K(y))x$. Maka kita mempunyai
\begin{align*}
D(M,N,\num 0) & \red \num 0 (K(N)) M \red M \text{ dan}\\
D(M,N,\num 1) & \red \num 1 (K(N)) M \red K(N) M \red N,
\end{align*}
seperti yang diperlukan.
\end{proof}

Gagasannya ialah bahwa $D(M,N)$ merepresentasikan pasangan $\tuple{M,N}$;
jika $P$ dianggap merepresentasikan pasangan semacam itu, maka $P(\num 0)$ dan
$P(\num 1)$ merepresentasikan proyeksi kiri dan kanan, $(P)_0$ dan $(P)_1$.
Kita akan menggunakan notasi terakhir ini.

\begin{lem}
Fungsi-fungsi yang !!{lambda definable} tertutup terhadap rekursi primitif.
\end{lem}

\begin{proof}
Kita perlu menunjukkan bahwa, jika diberikan sebarang suku $G$ dan $H$, kita
dapat menemukan suatu suku $F$ sedemikian sehingga
\begin{align*}
F(\num{0}) & \equiv G \\
F(\num{n+1}) & \equiv H(\num{n}, F(\num{n}))
\end{align*}
untuk setiap bilangan asli~$n$. Gagasan kasarnya ialah menghitung barisan
\emph{pasangan}
\[
\tuple{\num 0, F(\num 0)}, \tuple{\num 1, F(\num 1)}, \dots,
\]
dengan menggunakan numeral sebagai iterator. Perhatikan bahwa pasangan pertama
hanyalah $\tuple{\num 0,G}$. Jika diberikan pasangan
$\tuple{\num{n},F(\num{n})}$, pasangan berikutnya,
$\tuple{\num{n+1},F(\num{n+1})}$, seharusnya ekuivalen dengan
$\tuple{\num{n+1},H(\num{n},F(\num{n}))}$. Kita akan merancang suku
lambda~$T$ yang menjalankan transisi satu-langkah ini.

Perinciannya sebagai berikut. Definisikan $T(u)$ dengan
\[
T(u) = \tuple{S((u)_0), H((u)_0,(u)_1)}.
\]
Sekarang mudah diperiksa bahwa untuk setiap bilangan~$n$,
\[
T(\tuple{\num{n}, M}) \red \tuple{\num{n+1}, H(\num{n}, M)}.
\]
Seperti disarankan di atas, jika $G$ dan $H$ diberikan, definisikan~$F(u)$
dengan
\[
F(u) = (u(T,\tuple{\num 0, G}))_1.
\]
Dengan kata lain, pada masukan~$\num{n}$, $F$ mengiterasikan $T$ sebanyak $n$
kali mulai dari $\tuple{\num 0,G}$, kemudian mengembalikan komponen kedua.
Untuk langkah awal, kita mempunyai
\begin{enumerate}
\item $\num{0} (T, \tuple{\num 0, G}) \equiv \tuple{\num{0}, G}$
\item $F(\num{0}) \equiv G$
\end{enumerate}
Dengan induksi pada~$n$, dapat ditunjukkan bahwa bagi setiap bilangan asli
berlaku hal-hal berikut:
\begin{enumerate}
\item $\num{n+1}(T, \tuple{\num{0}, G}) \equiv \tuple{\num{n+1}, F(\num{n+1})}$
\item $F(\num{n+1}) \equiv H(\num{n}, F(\num{n}))$
\end{enumerate}
Untuk klausa kedua, kita mempunyai
\begin{align*}
F(\num{n+1}) & \red (\num{n+1}(T, \tuple{\num 0, G}))_1 \\
& \equiv (T(\num{n} (T, \tuple{\num 0, G})))_1 \\
& \equiv (T(\tuple{\num{n}, F(\num{n})}))_1 \\
& \equiv (\tuple{\num{n+1}, H(\num{n}, F(\num{n}))})_1 \\
& \equiv H(\num{n}, F(\num{n})).
\end{align*}
Di sini kita telah menggunakan hipotesis induksi pada baris kedua dari bawah.
Untuk klausa pertama, kita mempunyai
\begin{align*}
\num{n+1} (T, \tuple{\num 0, G}) &
\equiv T(\num{n} (T, \tuple{\num 0, G})) \\
& \equiv T( \tuple{\num{n}, F(\num{n})}) \\
& \equiv \tuple{\num{n+1}, H(\num{n}, F(\num{n}))} \\
& \equiv \tuple{\num{n+1}, F(\num{n+1})}.
\end{align*}
Di sini kita menggunakan klausa kedua pada baris terakhir. Jadi, kita telah
menunjukkan $F(\num 0)\equiv G$ dan, untuk setiap~$n$,
$F(\num {n+1})\equiv H(\num n,F(\num{n}))$, tepat seperti yang diperlukan.
\end{proof}

\end{document}

